\corrigesChapitre{2}{L’air, les gaz et la pression}
\par\noindent\begin{minipage}{\linewidth}
\begin{corrige}{1}{Les deux principaux gaz}
\begin{questions}
\item Le diazote et le dioxygène.
\item Environ 78 \% et 21 \% en volume ; les autres gaz représentent environ 1 \%.
\end{questions}
\end{corrige}
\end{minipage}
\par\noindent\begin{minipage}{\linewidth}
\begin{corrige}{2}{Air ou dioxygène ?}
\begin{questions}
\item Non : il contient plusieurs espèces chimiques.
\item Non : le dioxygène est un de ses constituants.
\end{questions}
\end{corrige}
\end{minipage}
\par\noindent\begin{minipage}{\linewidth}
\begin{corrige}{3}{Compresser}
\begin{questions}
\item C’est une compression.
\item Non si la seringue reste étanche : les molécules sont plus rapprochées.
\end{questions}
\pcFigureEx{figures/cor-compression.pdf}{Le récipient est étanche : le nombre de particules reste le même.}
\end{corrige}
\end{minipage}
\par\noindent\begin{minipage}{\linewidth}
\begin{corrige}{4}{Des unités de pression}
\begin{questions}
\item $1\,020\times100=102\,000\unite{Pa}$.
\item $2\times100\,000=200\,000\unite{Pa}$.
\end{questions}
\end{corrige}
\end{minipage}
\par\noindent\begin{minipage}{\linewidth}
\begin{corrige}{5}{Dans cent millilitres}
\begin{questions}
\item $21\unite{mL}$.
\item Environ $1\unite{mL}$, avec les arrondis du modèle.
\end{questions}
\end{corrige}
\end{minipage}
\par\noindent\begin{minipage}{\linewidth}
\begin{corrige}{6}{Un prélèvement plus grand}
\begin{questions}
\item $0{,}21\times2=0{,}42\unite L=420\unite{mL}$.
\item $0{,}78\times2=1{,}56\unite L$.
\end{questions}
\par\smallskip\noindent\textbf{Pour comprendre.}\enspace Une proportion de 21 \% signifie $21/100$ du volume total : $2\times21/100=0{,}42\unite L$. Pour le diazote : $2\times78/100=1{,}56\unite L$.\par
\end{corrige}
\end{minipage}
\par\noindent\begin{minipage}{\linewidth}
\begin{corrige}{7}{Lire un capteur}
\begin{questions}
\item Oui, de $0{,}4\unite{bar}$.
\item Non : si le gaz est enfermé, sa masse reste constante. Pression et masse sont des grandeurs différentes.
\end{questions}
\end{corrige}
\end{minipage}
\par\noindent\begin{minipage}{\linewidth}
\begin{corrige}{8}{Une masse d’air}
\begin{questions}
\item $m=1{,}2\times5=6{,}0\unite g$.
\item La température et la pression ; la valeur donnée ne convient pas à toutes les conditions.
\end{questions}
\end{corrige}
\end{minipage}
\par\noindent\begin{minipage}{\linewidth}
\begin{corrige}{9}{Peser par différence}
\begin{questions}
\item $425{,}6-423{,}2=2{,}4\unite g$.
\item $V=2{,}4/1{,}2=2{,}0\unite L$, si les conditions justifient cette masse volumique.
\end{questions}
\par\smallskip\noindent\textbf{Pour comprendre.}\enspace Soustraire les masses du même ballon : $425{,}6-423{,}2=2{,}4\unite g$. Puis $V=m/\rho=2{,}4/1{,}2=2\unite L$, avec les conditions de température et de pression annoncées.\par
\end{corrige}
\end{minipage}
\par\noindent\begin{minipage}{\linewidth}
\begin{corrige}{10}{Modéliser une détente}
\begin{questions}
\item Agrandir l’espace du récipient et répartir davantage les particules.
\item Le même nombre et les mêmes types de particules : aucun constituant n’est créé.
\end{questions}
\pcFigureEx{figures/cor-detente.pdf}{Même nombre et mêmes types de particules ; elles se répartissent dans un volume plus grand.}
\end{corrige}
\end{minipage}
\par\noindent\begin{minipage}{\linewidth}
\begin{corrige}{11}{Une conclusion excessive}
\begin{questions}
\item Non. Une seule observation qualitative ne prouve pas une loi valable dans toutes les situations.
\item La température, mais aussi la quantité de gaz et l’étanchéité. Le cours n’impose ici qu’une évolution qualitative.
\end{questions}
\end{corrige}
\end{minipage}
\par\noindent\begin{minipage}{\linewidth}
\begin{corrige}{12}{L’air humide}
\begin{questions}
\item Parce que de la vapeur d’eau s’ajoute au mélange et modifie les fractions de l’ensemble.
\item Non : les proportions de toutes les espèces doivent totaliser 100 \%. Les 78 \% et 21 \% se rapportent à l’air sec.
\end{questions}
\end{corrige}
\end{minipage}
\par\noindent\begin{minipage}{\linewidth}
\begin{corrige}{13}{Chercher une fuite}
\begin{questions}
\item Une fuite peut réduire la quantité de gaz ; le gaz peut aussi se refroidir après une compression rapide.
\item Attendre l’équilibre thermique, contrôler les raccords et répéter l’essai. Une baisse seule ne permet pas de choisir.
\end{questions}
\end{corrige}
\end{minipage}
\par\noindent\begin{minipage}{\linewidth}
\begin{corrige}{14}{Doubler l’échantillon}
\begin{questions}
\item $52{,}5\unite{mL}$ puis $105\unite{mL}$.
\item Non : les deux quantités ont doublé, mais le pourcentage reste $21\,\%$.
\end{questions}
\end{corrige}
\end{minipage}
\par\noindent\begin{minipage}{\linewidth}
\begin{corrige}{15}{Un modèle incomplet}
\begin{questions}
\item L’air est un mélange et un gaz occupe tout le volume ; ses particules sont en mouvement.
\item Utiliser plusieurs codes pour les espèces, répartis dans toute la boîte. Signaler que positions et tailles sont schématiques.
\end{questions}
\pcFigureEx{figures/cor-air.pdf}{Schéma qualitatif : un gaz occupe tout le récipient et ses particules sont en mouvement ; les couleurs codent des espèces.}
\end{corrige}
\end{minipage}
\par\noindent\begin{minipage}{\linewidth}
\phantomsection\label{cor-defi-2}\pcSolutionDefi{2}{Dimensionner un modèle}\begin{questions}
\item $156$ pour le diazote, $42$ pour le dioxygène et $2$ pour les autres gaz.
\item $1\,000/200=5\unite{mL}$ par jeton, en proportion de volume et dans les mêmes conditions.
\item Le modèle ne montre ni les distances réelles ni les mouvements. Un jeton ne correspond pas à une molécule réelle de cinq millilitres.
\end{questions}
\end{minipage}
\par\noindent\begin{minipage}{\linewidth}
\phantomsection\label{cor-activite-2}\pcSousTitre{Pistes pour l’investigation}
La pression augmente lorsque le volume diminue lentement, à quantité de gaz conservée et après retour à une température stable. Le nombre de particules n’augmente pas. Le retour au volume initial doit donner une pression voisine de la première valeur ; sinon, rechercher fuite, température ou lecture. Cette expérience qualitative ne prouve pas à elle seule une loi quantitative universelle.\par\end{minipage}
